\documentclass[11pt,a4paper]{article}
\usepackage[utf8]{inputenc}
\usepackage[T1]{fontenc}
\usepackage{amsmath,amssymb}
\usepackage{booktabs}
\usepackage{listings}
\usepackage{xcolor}
\usepackage{hyperref}
\usepackage{geometry}
\geometry{margin=2.5cm}

\lstset{
  basicstyle=\ttfamily\small,
  breaklines=true,
  frame=single,
  backgroundcolor=\color{gray!10},
  numbers=left,
  numberstyle=\tiny
}

\title{02244 - Logic for Security\\
\vspace{0.5em}
Project on Security Protocols}

\author{%
  \textbf{Authors}\\[0.5em]
  Zeinab\\
  Fahmida\\
  Mirtha\\[1em]
  March 2026
}

\date{}

\begin{document}
\maketitle
\thispagestyle{empty}

\newpage
\tableofcontents
\newpage

\section{Introduction - Zeinab}

\noindent
This project develops and verifies an Open Authorization protocol for the following setting.
User $A$ stores personal photos at server $B$ and wants to use a photo-book service $P$.
Instead of downloading from $B$ and uploading to $P$, user $A$ should be able to delegate
restricted access to $P$ directly at $B$ (read-only, scoped access).

The trust architecture includes an identity provider $\mathit{IDP}$.
User $A$ has no long-term public/private key pair in our core design assumptions and authenticates through a shared password-like secret with $\mathit{IDP}$.
This secret is modeled as low-entropy and is never used as an encryption key.
All other principal keys and authorization statements are represented in symbolic AnB/OFMC models.

The report follows the required week-by-week development methodology:
baseline modeling under Dolev-Yao, lazy intruder analysis, typing/format hardening,
channel abstraction, and guessable-password verification.
An important design decision throughout is to preserve executable models even when attacks are found,
because those attack traces are used to motivate and justify each refinement step.

\newpage
\section{Open Authorization Protocol Development (Week 2--3) - Zeinab}

\subsection{Initial baseline (Week 2)}

The first executable protocol model is intentionally simple:
\begin{enumerate}
  \item $A \rightarrow P$: signed delegation context $(A,B,P,F,T)$.
  \item $P \rightarrow B$: forwards signed context and requests folder $F$.
  \item $B \rightarrow P$: returns encrypted photos.
\end{enumerate}

The Week 2 objective is not final security, but establishment of a minimal
formal baseline that can be analyzed under Dolev-Yao deductions.
The static analysis confirms a confidentiality weakness in this baseline:
if the delegate role is intruder-controlled, the photo payload may be derivable.
This gives a concrete reason to move to identity-provider-based delegation in Week 3.

\subsection{Identity-provider-based delegation (Week 3)}

In Week 3 we redesign the flow so that authorization is issued by $\mathit{IDP}$:
\begin{enumerate}
  \item $P \rightarrow A$: access request with scope and freshness.
  \item $A \rightarrow \mathit{IDP}$: request plus password proof $h(pw(A,\mathit{IDP}),N)$.
  \item $\mathit{IDP} \rightarrow A$: signed delegation grant.
  \item $A \rightarrow P$: forward signed grant.
  \item $P \rightarrow B$: present grant and requested scope.
  \item $B \rightarrow P$: encrypted photo data.
\end{enumerate}

This satisfies the assignment's modeling requirement that the password is used
for authentication evidence only and not as a transport encryption key.
However, OFMC still finds attacks in this stage (secrecy and/or weak authentication depending on variant),
which is expected in an iterative design process.

\subsection{Special task (Week 3): lazy intruder}

The role-executability task was solved explicitly using the lecture method:
with $P=i$ and honest $A,B,\mathit{IDP}$, each outgoing message of role $P$
is derived from prior knowledge using Axiom/Compose plus lazy variables.
This proves role executability independently from whether global security goals hold.

\subsection{Representative AnB model}
\begin{lstlisting}
Protocol: OpenAuth_W3_DelegationViaIDP

Types:
  Agent A,B,P,I;
  Number F,T,NA,Grant,Photos;
  Function pk,pw,h;

Actions:
  P -> A: P,B,F,T,NA
  A -> I: A,P,B,F,T,NA,h(pw(A,I),NA)
  I -> A: {A,P,B,F,T,NA,Grant}inv(pk(I))
  A -> P: {A,P,B,F,T,NA,Grant}inv(pk(I))
  P -> B: {A,P,B,F,T,NA,Grant}inv(pk(I)),F
  B -> P: {Photos}pk(P)
\end{lstlisting}

\newpage
\section{Typing, Formats, and Key Lookup (Week 4) - Fahmida}

\subsection{Motivation}

Week 4 requires two linked improvements:
\begin{itemize}
  \item Replace ambiguous message concatenations with explicit formats.
  \item Remove the assumption that $A$ initially knows all public keys by adding a separate key-lookup protocol.
\end{itemize}

The key engineering idea is to make message intent explicit in syntax.
Different logical message purposes should have different top-level constructors
so that unification across incompatible message kinds becomes impossible or controlled.

\subsection{Main protocol with formats}

The main delegation model introduces constructors such as
\texttt{fReq}, \texttt{fAuthProof}, \texttt{fGrant}, \texttt{fUse}, and \texttt{fData}.
During development, several over-wrapped variants became non-executable in OFMC.
The final Week 4 file is executable, but still admits weak-auth attacks.
This is acceptable at this stage because the Week 4 goal is primarily typing discipline,
not full protocol closure.

\subsection{GetPublicKey protocol}

A separate protocol was modeled so that $A$ can query $\mathit{IDP}$ for a target key:
\begin{enumerate}
  \item $A \rightarrow \mathit{IDP}$: request key for principal $X$.
  \item $\mathit{IDP} \rightarrow A$: signed response containing $(X,pk(X))$.
\end{enumerate}

This decomposition follows the assignment requirement and improves realism of key-distribution assumptions.

\subsection{Special task (Week 4): type-flaw resistance argument}

The argument is documented in text form: if two non-variable message patterns unify,
they must share compatible top-level format structure.
In this model family, message kinds are separated by dedicated constructors,
which significantly reduces type-confusion opportunities.

\subsection{Representative AnB model}
\begin{lstlisting}
Protocol: OpenAuth_W4_KeyLookup

Types:
  Agent A,I,X;
  Function pk,fKeyReq,fKeyResp;

Actions:
  A -> I: fKeyReq(A,X)
  I -> A: {fKeyResp(X,pk(X))}inv(pk(I))

Goals:
  A authenticates I on fKeyResp(X,pk(X));
\end{lstlisting}

\newpage
\section{Channels, Guessable Passwords, and Final Status (Week 5--6) - Mirtha}

\subsection{Week 5: pseudonymous channel abstraction}

We modeled a TLS-style abstraction using OFMC pseudonymous channel notation.
The objective was to replace explicit low-level transport cryptography with secure-channel assumptions while preserving business-level authorization semantics.
The Week 5 model is executable, but OFMC reports an attack on secrecy.
This result shows that channel abstraction alone does not eliminate authorization design weaknesses.

\subsection{Week 6: guessable password verification}

Week 6 focuses on low-entropy credential robustness.
Three checks were performed:
\begin{enumerate}
  \item Baseline guessable-password model: OFMC reports \texttt{ATTACK\_FOUND} with goal \texttt{guesswhat}.
  \item Hardened model with additional fresh value in password proof: no immediate \texttt{guesswhat}, but other secrecy attacks remain.
  \item Focused guessable-only model: bounded check (\texttt{--numSess 1}) reports \texttt{NO\_ATTACK\_FOUND}.
\end{enumerate}

Therefore, for the dedicated Week 6 special-task goal and tested session bound,
the hardened/focused construction provides evidence of guessing resistance.
At the same time, remaining secrecy/authentication issues in broader models are explicitly acknowledged.

\subsection{Final protocol status}

The project outcome is not ``one perfect protocol,'' but a justified development path:
\begin{itemize}
  \item Week 2 establishes a minimal baseline and exposes initial leakage.
  \item Week 3 introduces identity-provider trust and completes lazy-intruder executability.
  \item Week 4 adds typed formats and key-lookup decomposition.
  \item Week 5 explores channel-level abstraction.
  \item Week 6 isolates and verifies guessable-password behavior under bounded analysis.
\end{itemize}

This progression is aligned with the assignment's expectation that report and logbook document reasoning, attacks, and refinements---not only final syntax.

\subsection{Summary table of special tasks}
\begin{center}
\begin{tabular}{ll}
\toprule
\textbf{Week} & \textbf{Special Task Outcome} \\
\midrule
2 & Static Dolev-Yao analysis completed (baseline weakness identified). \\
3 & Lazy intruder executability of role $P$ proved step-by-step. \\
4 & Type-flaw resistance argument and key-lookup protocol provided. \\
5 & Pseudonymous-channel variant implemented and analyzed. \\
6 & Guessable-password check completed (bounded no-attack in focused model). \\
\bottomrule
\end{tabular}
\end{center}

\newpage
\section*{Delivered AnB and support files}
\begin{itemize}
  \item \texttt{Week2\_Problem2\_First\_AnB\_Version.anb}
  \item \texttt{Week3\_Problem1\_AnB\_With\_IdentityProvider.anb}
  \item \texttt{Week3\_Problem3\_Refined\_HonestProvider.anb}
  \item \texttt{Week4\_Problem1\_MainProtocol\_Formats.anb}
  \item \texttt{Week4\_Problem2\_AnB\_KeyLookup\_Protocol.anb}
  \item \texttt{Week5\_Problem1\_PseudonymousChannels.anb}
  \item \texttt{Week6\_Problem1\_GuessablePassword\_Check.anb}
  \item \texttt{Week6\_Problem2\_GuessablePassword\_Hardened.anb}
  \item \texttt{Week6\_Problem3\_GuessablePassword\_OnlyCheck.anb}
  \item Supporting notes: Week2/3/4/5/6 markdown analysis files.
\end{itemize}

\section*{References}
\begin{itemize}
  \item 02244 course slides and assignment handout.
  \item OFMC 2024 and AnB modeling documentation.
  \item Week-specific protocol analysis notes created during this project.
\end{itemize}

\end{document}
